% SOURCES: docs/0.2.0v/plan.md (§4 Guiding Principles, §6 Contracts, §17 Milestones,
%          §7 Day 0); CLAUDE.md (Sprint Model, Git Workflow); .claude/rules/*;
%          docs/1.0.0v/supervisor-report.md (§8 Build & CI)
% GUIDANCE: Describe HOW the team worked, not what was built. This is a strong
% chapter for this project because the process is unusually well-documented.

\chapter{Work Methodology}
\label{ch:methodology}

\section{Development Model: Task-Driven Milestones}
% DRAFT (review):
Development followed a \emph{task-driven, milestone-gated} model rather than a
fixed time-box. Progress was unlocked by completing the tasks that gate each
milestone, from \textbf{M1 (Foundation)} through \textbf{M5 (Polish + Demo)} and
the \textbf{v1.0.0} demo cut. Every task carried a single owner, a priority
(\texttt{P0} blocks the demo, \texttt{P1} is a Tier-1 feature, \texttt{P2} is
polish), and explicit dependency links.
% TODO: cite the milestone table (reproduce a compact version here or in ch. 4).

\section{Team Structure and Ownership}
% SOURCES: CLAUDE.md (Module Ownership); docs/0.2.0v/plan.md (§10 Tasks)
% DRAFT (review):
Three engineers owned disjoint areas to avoid merge contention:
\begin{itemize}[leftmargin=1.5em]
  \item \textbf{Document Model, generator, runtime, SEO, export, shell, CI} ---
        the output pipeline and Electron shell.
  \item \textbf{State stores} --- Zustand stores, undo/redo history via
        \texttt{immer} patches, persistence, autosave/crash-recovery.
  \item \textbf{User interface} --- canvas, panels, sidebar, topbar, layers tree.
\end{itemize}
% TODO: confirm names/IDs and add a one-line responsibility per member.

\section{Cross-Owner Contracts}
% SOURCES: docs/0.2.0v/plan.md (§6 Contracts C1–C12); docs/0.2.0v/architecture.md (§6)
% DRAFT (review):
Interfaces between owners were fixed as twelve numbered \emph{contracts}
(C1--C12): the shared types, the Zod schemas, the operation set, the typed IPC
surface, the store hooks, the generator entry point, the presets registry, the
validation function, the token resolver, the semantic-inference hint, the image
IPC contract, and the export orchestrator. Any breaking change to a contract
required a labelled pull request and review from the named downstream consumer,
which kept parallel work from silently diverging.

\section{Quality Gates and Tooling}
% SOURCES: .claude/rules/git.md, preflight.md; CLAUDE.md (Commands, Code Quality);
%          docs/1.0.0v/supervisor-report.md (§8)
% DRAFT (review):
Every push ran three gates --- lint (ESLint + Prettier), type-checking (strict
TypeScript across all \texttt{tsconfig}s), and the test suite --- enforced locally
by a pre-push hook and again in continuous integration. Specialised review passes
checked output determinism, the no-\texttt{position:absolute} invariant,
accessibility, and contract changes. A weekly integration ritual merged open
branches and ran the full matrix.
% TODO: note the ``no AI attribution'' policy if relevant to your academic
% integrity statement.

\section{Version Control and Release Workflow}
% SOURCES: .claude/rules/git.md; .github/workflows/{ci,release}.yml; CHANGELOG.md
% DRAFT (review):
Work used conventional commits on short-lived feature branches squash-merged into
\texttt{main}. Releases were tag-driven: pushing a \texttt{v*} tag triggered a
per-operating-system build matrix that produced native installers and published a
GitHub Release. Versions progressed v0.1.0 → v0.2.0 → v0.3.0 → \textbf{v1.0.0}.
